\documentclass[11pt, a4paper]{report}
% --- Core Mathematics Packages ---
\usepackage{amsmath, amssymb, amsthm}
\usepackage{mathtools} % For advanced formatting (e.g., \coloneqq for :=)

\usepackage{fancyhdr}
%\setlength{\headheight}{13.59999pt}
%\addtolength{\topmargin}{-1.59999pt}

% --- Typography and Page Layout ---
%\usepackage[margin=1in]{geometry}

\usepackage[a4paper, top=1in, bottom=1in, left=1.4in, right=1.4in]{geometry}

\usepackage{microtype} % Improves text justification and spacing
\usepackage{xcolor}    % Required for coloring links
\usepackage{hyperref}  % Makes the table of contents and cross-references clickable
\hypersetup{
    colorlinks=true,
    linkcolor=black,
    filecolor=magenta,      
    urlcolor=cyan,
    pdftitle={Mathematics Notes},
}
\usepackage{nameref} % Readable theorem referencing

\usepackage{subfiles}

\usepackage{framed}
\definecolor{shadecolor}{gray}{0.92}

\usepackage{tikz}
\usetikzlibrary{arrows.meta}
\usetikzlibrary{decorations.pathreplacing,calligraphy}


% --- Theorem Environments ---
% This links the numbering of theorems, lemmas, and definitions together.
\theoremstyle{definition}
\newtheorem{definition}{Definition}[chapter]
\newtheorem{example}[definition]{Example}

\theoremstyle{plain}
\newtheorem{theorem}[definition]{Theorem}
\newtheorem{lemma}[definition]{Lemma}
\newtheorem{corollary}[definition]{Corollary}
\newtheorem{proposition}[definition]{Proposition}

\theoremstyle{remark}
\newtheorem*{remark}{Remark} % The asterisk removes numbering for remarks


\newtheorem{exercise}{Exercise}[chapter]
\newtheorem*{exercise*}{Exercise}  % unnumbered, for one-off inline prompts


% --- Custom Shortcuts (Macros) ---
% Blackboard bold letters for standard sets
\newcommand{\R}{\mathbb{R}}
\newcommand{\N}{\mathbb{N}}
\newcommand{\Z}{\mathbb{Z}}
\newcommand{\Q}{\mathbb{Q}}
\newcommand{\C}{\mathbb{C}}

% Upright 'd' for differentials (Standard UK/Imperial convention)
\newcommand{\dd}{\mathrm{d}}
%%% TITLE PAGE
\ifx\nauthor\undefined
  \def\nauthor{Your Name}
\fi

\author{Based on lectures by \nlecturer \\ \small Notes taken by \nauthor}

\date{
  \nterm\ \nyear \\
  \small Compiled on \today
  \endgraf
  \bigskip
  \parbox{0.9\textwidth}{\centering
    These notes are not in any way official and do not represent
    the content of lectures. Be mindful of any mistakes or omissions
    if using these notes for self-study or revision.
    Any and all errors are my own.
  }
  \endgraf
  \medskip
  More details at \href{https://stones-y.github.io}{https://stones-y.github.io}.
}

\title{\nmodule}

\AtBeginDocument{
  \hypersetup{
    pdfauthor={\nauthor},
    pdftitle={\nmodule},
  }
}

%%% HEADERS
\pagestyle{fancy}
\fancyhf{}
\fancyhead[L]{\nouppercase{\leftmark}}
\fancyhead[R]{\small\nmodule}
\fancyfoot[C]{\thepage}





\def\nmodule{MATH40002 - Analysis I}
\def\nlecturer{Dr. M.-A. Lawn and Mr M. Mendez Guaraco}
\def\nterm{Autumn 2025 -- Spring 2026}
\def\nyear{}

\begin{document}
\maketitle
\tableofcontents
\newpage

\chapter{Numbers and Sets}
\section{Introduction}
In the introductory module (MATH40001), we constructed the number systems $\N,\Z,\Q$ and $\R$.
Real analysis focuses on the real numbers, $\R$, formalising many properties about functions that we know from calculus.
For example, an intuitive notion of a continuous function is that you can draw the function on a piece of paper
without taking your pen off the page. We make concepts like these precise in this module.

\section{Countability}
Before starting this chapter (and, in general, this module), it
is imperative that you are familiar with the logic and sets part
of the introductory module, MATH40001.
\begin{definition}
    \begin{shaded*}
        A set $A$ is countably infinite if there exists a bijection $f:\N\to A$.
    \end{shaded*}
\end{definition}
\noindent Intuitively, this says that you can list the elements of $A$ as $a_1,a_2,\dots$ using only the natural numbers.
We say that a set is \textit{countable} if it is finite or countably infinite. Note we will take $\N=\{1,2,\dots\}$, unless stated otherwise.\par
Note that if there is a bijection between a given pair of sets $A$ and $B$, then $A$ is countable if and only if $B$ is countable.
This extends the natural idea of finite sets having the same number of elements if and only if there is a bijection between the sets. Formally, we define two sets having the same $cardinality$ as follows:
\begin{definition}
    \begin{shaded*}
        Let $A,B$ be sets. We say that $A$ and $B$ have the same cardinality if there exists a bijection $f:A\to B$.
    \end{shaded*}
\end{definition}
The idea is that even in infinite settings, if we know a property about the 'size' of one set $A$, and there exists a bijection with another set $B$, then $B$ shares this notion of 'size'.
For example, as we mentioned, if such a bijection exists between two sets $A,B$, then $A$ is countable if and only if $B$ is countable. 
This gives rise to an equivalence relation, which we define as cardinality.
\begin{exercise*}
    Show that the relation $\sim$ on $\mathcal{P}(\mathcal{U})$ (where $\mathcal{U}$ is the universal set of interest) defined by $A\sim B \iff \exists f:A\to B$ such that $f$ is a bijection, is an equivalence relation.
\end{exercise*}


\begin{theorem}
    \begin{shaded*}
        If $A\subseteq\N$ is an infinite set, then $A$ is countably infinite.
    \end{shaded*}
\end{theorem}
\begin{proof}
    We construct a function $f:\mathbb{N}\to A$ inductively: since $A$ is nonempty, and $A\subseteq\N$,
    by well-ordering of the natural numbers we can pick $a_1=\min A$, and let $f(1)=a_1$. Assume that, for $k\ge 1$, $a_1,\dots,a_k$
    have been chosen and $f(i)=a_i$ for $1\le i \le k$. Since $A$ is
    infinite, $A\setminus\{a_1,\dots,a_k\}$ is nonempty. Pick $a_{k+1}=\min\left(A\setminus\{a_1,\dots,a_k\}\right)$,
    and let $f(k+1)=a_{k+1}$. Thus by construction $f$ is strictly increasing. We claim that $f$ is a bijection.
    \begin{itemize}
        \item Injective: this holds because by construction $f$ is strictly increasing. We saw in the introductory module that strictly increasing functions are injective.
        \item Surjective: suppose for contradiction that $f$ is not surjective. Then there exists
        $a\in A$ such that for all $n\in\N$, $f(n)\ne a$. Consider $A\setminus\{\}$.
    \end{itemize}
\end{proof}

\begin{lemma}
    \begin{shaded*}
        Let $S$ be an infinite set. There exists an injection $f:\N\to S$.
    \end{shaded*}
\end{lemma}
\begin{proof}
    We construct a function $f:\N\to S$ inductively: since
    $S$ is infinite, it is nonempty, so we can pick $s_1\in S$ and define $f(1)\coloneqq s_1$.
    Suppose for $k\in\N$, $s_1,\dots,s_k$ have been chosen and we have defined $f(i)=s_i$ for $1\le i\le k$. Since
    $S$ is infinite, $S\setminus\{s_1,\dots,s_k\}$ is nonempty, so there is $s_{k+1}\in S\setminus \{s_1,\dots,s_k\}$. Define $f(k+1)\coloneqq s_{k+1}$.
    By construction, if $i<j$ are natural numbers, then $s_j\in S\setminus\{s_1,\dots,s_{j-1}\}$ so $s_j\ne s_i$. So $f$ is an injection.\qedhere
\end{proof}
As a result, we can always $embed$ the natural numbers into any infinite set via an injection $f:\N\to S$ (where $S$ is any infinite set).
In this way, the natural numbers can be thought of as 'the smallest type of infinity'.

A not so surprising fact which we can prove immediately is that the set of integers, $\Z$, is countably infinite:
\begin{theorem}
    \begin{shaded*}
        $\Z$ is countably infinite.
    \end{shaded*}
\end{theorem}
\begin{proof}
    Define $f:\N\to \Z$ by $$f(n)=\begin{cases}
    n/2, & \text{if } n\text{ even}\\
    -\frac{(n+1)}{2}, & \text{if }n\text{ odd}
    \end{cases}$$
    Then $f$ is injective: suppose $f(a)=f(b)$. If $a$ and $b$ are of different parity then $f(a)$ and $f(b)$ are of different sign, so cannot be equal. So this case cannot happen.
    If $a$ and $b$ are even, then $a/2=b/2$ so $a=b$. If $a$ and $b$ are odd, then $-(a+1)/2=-(b+1)/2$ so $a=b$. So $f$ is injective.
    Also, $f$ is surjective: let $m\in\Z$ be arbitrary. If $m<0$ then $f(-2m-1)=m$, and if $m\ge 0$ then $f(2m)=m$. Thus $f$ is a bijection and $\Z$ is countable.\qedhere
\end{proof}
\begin{remark}
    This proof used a useful trick called interleaving, where we listed the integers as $0,-1,1,-2,2,\dots$, interleaving the negative and positive integers to create the desired bijection between $\N$ and $\Z$.
\end{remark}  
A less trivial fact is that the rational numbers, $\Q$, is countable. This may be hard to see until observing the proof of the following fact:
\begin{theorem}
\label{thm:N2}
    \begin{shaded*}
        $\N^2$ is countably infinite.
    \end{shaded*}
\end{theorem}
\begin{proof}
    \begin{figure}[h]
    \centering
    \begin{tikzpicture}[>={Stealth[length=4pt]}, scale=1.2]

        % Dots at each grid point
        \foreach \i in {1,...,5} {
          \foreach \j in {1,...,5} {
            \fill[blue!70] (\i,-\j) circle (2pt);
          }
        }
        
        % Diagonal path arrows
        \draw[->] (1,-1) -- (2,-1);
        \draw[->] (2,-1) -- (1,-2);
        \draw[->] (1,-2) -- (1,-3);
        \draw[->] (1,-3) -- (2,-2);
        \draw[->] (2,-2) -- (3,-1);
        \draw[->] (3,-1) -- (4,-1);
        \draw[->] (4,-1) -- (3,-2);
        \draw[->] (3,-2) -- (2,-3);
        \draw[->] (2,-3) -- (1,-4);
        \draw[->] (1,-4) -- (1,-5);
        \draw[->] (1,-5) -- (2,-4);
        \draw[->] (2,-4) -- (3,-3);
        \draw[->] (3,-3) -- (4,-2);
        \draw[->] (4,-2) -- (5,-1);
        
        % Continuation dots
        \foreach \i in {1,...,5} {
          \node at (\i, -5.6) {$\vdots$};
        }
        \foreach \j in {1,...,5} {
          \node at (5.6, -\j) {$\cdots$};
        }
        
        \end{tikzpicture}
        \caption{The diagonal listing of $\mathbb{N}^2$, showing it is countably infinite.}
        \label{fig:diagonal-listing}
    \end{figure}
    The idea is illustrated in the above diagram. Note that the listing goes $$(0,0),(0,1),(1,0),(2,0),(1,1),(0,2),\dots$$
    Notice that along each diagonal, $a+b$ is constant for all pairs $(a,b)$ in that diagonal. This, along with the triangular number formula,
    provides the motivation for how to construct a bijection between $\N^2$ and $\N$.\par
    Explicitly, define $f:\N^2\to\N$ by $f(a,b)=\frac{(a+b)(a+b+1)}{2}+b$. This is a bijection (exercise) so $\N^2$ is countable. Note that here we took $\N=\{0,1,2,\dots\}$.\qedhere
\end{proof}
As a consequence, we get:
\begin{theorem}
\label{thm:qcount}
    \begin{shaded*}
        $\Q$ is countably infinite.
    \end{shaded*}
\end{theorem}
\begin{proof}
    First we use Theorem \ref{thm:N2} to show that $\Q_{\ge 0}\coloneqq \{x\in\Q:x\ge 0\}$ is countably infinte.
    Let $$S=\{(a,b)\in\N^2: b\ge 1,\, \gcd(a,b)=1\}.$$
    Then $f:S\to \Q_{\ge 0 }$, defined by $f(a,b)=a/b$ is a bijection because we know that all rational numbers in $\Q_{\ge 0}$ take the unique form $a/b$ where $b\ge 1$ and $\gcd(a,b)=1$.
    Since $\N^2$ is countably infinite by Theorem \ref{thm:N2}, $S\subseteq \N^2$ is also countably infinite, so $\Q_{\ge 0}$ is countably infinite.\par
    It follows (exercise) that $\Q$ is countably infinite, using the same idea (interleaving) as the proof that $\Z$ is countably infinite.
\end{proof}
\begin{exercise*}
Finish the proof of Theorem \ref{thm:qcount}. (\textit{Hint: use the fact we proved above that $\Q_{\ge 0}$ is countably infinite, so there exists a bijection }$h:\N\to\Q_{\ge0}$).
\end{exercise*}

The following lemma is a useful way of proving that a given set $S$ is countably infinite without explicitly constructing a bijection $f:\N\to S$.

\begin{lemma}[A criterion for countability]
\label{lem:countabilitycriterion}
    \begin{shaded*}
        Let $S$ be an infinite set.
        \begin{enumerate}
            \item[1.] $S$ is countably infinite if and only if there exists an injection $\iota: S\to\N$.
            \item[2.] $S$ is countably infinite if and only if there exists a surjection $g:\N\to S$.
        \end{enumerate}
    \end{shaded*}
\end{lemma}
\begin{proof}
    a
\end{proof}
The idea for item (1) above is that if you can $embed$ your set into the natural numbers via an injection $\iota:S\to\N$, then it must be countably infinite, as $\iota(S)\subseteq\N$, and we have shown infinite subsets of $\N$ are countable,
so since $\iota$ is an injection, $S$ has the same cardinality as $\iota (S)$, (i.e. $S\sim \iota(S)$), so $S$ is countably infinite. The idea for item (2)
is similar: if you can use the natural numbers to list all the elements of your infinite set $S$, then even though this surjective map $g:\N\to S$ may not be injective, we have nonetheless listed all the elements of $S$ using only the natural numbers,
so $S$ must be countably infinite also.\par
A simple example using Lemma \ref{lem:countabilitycriterion} is that finite products of countable sets are countable (note that we already have the tools to prove this without Lemma \ref{lem:countabilitycriterion} by usin Theorem \ref{thm:N2}, which is a more elementary way of proving this fact via an explicit bijection).\par
Before looking at the example, try and think of properties of the integers that we can use to construct an injection $\iota:A\times B\to \N$.
\begin{example}
    If $A,B$ are countably infinite sets, then $A\times B$ is countably infinite.
\end{example}
\begin{proof}
    Recall from the introductory module that we define $$A\times B\coloneqq \{(a,b):a\in A,\, b\in B\}.$$
    Since $A$ and $B$ are countably infinite, let $f:\N\to A$, $g:\N\to B$ be bijections. Therefore, we can write $A=\{a_1,a_2,\dots\}$ and $B=\{b_1,b_2,\dots\}$.
    Define $\iota: A\times B\to \N$ by $\iota(a_i,b_j)=2^i3^j$. By the Fundamental Theorem of Arithmetic from the introductory module, $\iota$ is an injection because prime factorisations are unique,
    so by Lemma \ref{lem:countabilitycriterion}, $A\times B$ is countably infinite.\qedhere
\end{proof}
This trick of mapping each element of the set to a prime factorisation is useful to keep in mind, as it is a very neat way of proving countability for sets.

Not all sets are countably infinite:
\begin{definition}
    \begin{shaded*}
        An infinite set is called $uncountable$ if it is not countably infinite.
    \end{shaded*}
\end{definition}

\noindent For example, the set of real numbers $\R$ is uncountable. To show this, we first prove:

\begin{lemma}
\label{lem:01}
    \begin{shaded*}
        The set $(0,1)\coloneqq\{x\in\R:0<x<1\}$ is uncountable.
    \end{shaded*}
\end{lemma}
\begin{proof}
    The main technique that we use in arguing why a given set is uncountable is \textit{Cantor's diagonal argument}, which goes like this.
    Suppose, for contradiction, that $(0,1)$ is countably infinite. Then there exists a bijection $f:\N\to (0,1)$ (note here we will let $\N=\{1,2,\dots\}$ for aesthetics; the argument does not change).
    For each $n\in\N$, we write $f(n)=x_n=0.a_{1n}a_{2n}\dots$, writing each $x_n$ in decimal form, excluding tails of the digit $9$ (for example, $0.8999...=0.9$).
    We now get a listing of the numbers in $(0,1)$ which looks like:
    \begin{align*}
        f(1)&=0.a_{11}a_{12}a_{13}a_{14}\dots\\
        f(2)&=0.a_{21}a_{22}a_{23}a_{24}\dots\\
        f(3)&=0.a_{31}a_{32}a_{33}a_{34}\dots\\
        f(4)&=0.a_{41}a_{42}a_{43}a_{44}\dots\\
        \vdots
    \end{align*}
    The idea is that we look at the $diagonal$ elements, $a_{ii}$ for $i\ge 1$, and construct a number $x\in (0,1)$ such that for all $i\ge 1$,
    the $i^{\mathrm{th}}$ decimal place of $x$ differs from the $i^{\mathrm{th}}$ decimal place of $f(i)$, thereby constructing a number $x\in(0,1)$ that is not in the enumeration $f$ of the interval $(0,1)$, which is a contradiction to $f$ being surjective.
    We construct $x\in (0,1)$ as follows: for $j\ge 1$, let $b_{j}$ be the $j^{\mathrm{th}}$ decimal place of $x$, then we define
    $$b_j=\begin{cases}2 & \text{if }a_{jj}=1\\ 1 & \text{if }a_{jj}\ne 1\end{cases}$$
    Then by construction, $b_i \ne a_{ii}$ for each $i\ge 1$, so in particular (by uniqueness of decimal representations under our restriction to decimal expansions with no digit $9$ tails), we have $x\ne f(i)$ for all $i\ge 1$,
    so $x\in (0,1)$ is not in the image of $f$, a contradiction to $f$ being surjective onto $(0,1)$. Therefore, $(0,1)$ is uncountable.\qedhere
\end{proof}
\begin{theorem}
    \begin{shaded*}
        $\R$ is uncountable.
    \end{shaded*}
\end{theorem}
\begin{proof}
    Follows from Lemma \ref{lem:01} (exercise). \qedhere
\end{proof}

\begin{exercise*}
    Find a bijection $f:(0,1)\to \R$, and deduce that $\R$ is an uncountable set.
\end{exercise*}

\chapter{Sequences and Series}

\section{Sequences}

\begin{definition}
    \begin{shaded*}
        A real sequence is a function $f:\N\to\R$.
    \end{shaded*}
\end{definition}
In general, $\R$ can be replaced by a metric space $X$, for example the complex numbers $\C$ which we will see later. It is therefore important to distinguish $real$ sequences. For notation, we will write $(a_n)_{n=1}^{\infty}$ to denote a sequence $a_1,a_2,\dots$, and sometimes this will be shortened to $(a_n)$, or just $a_n$. 


\subsection{Convergence of sequences}
Intuitively, a real sequence converges to a limit $L\in\R$ if we can make the terms of our sequence, $a_n$, aribtrarily close to $L$, whenever we go far enough out into the sequence.
Clearly, this does not suffice as a mathematical definition, so instead we define convergence of a real sequence $(a_n)$ as follows.
\begin{definition}[Convergence of a sequence in $\R$]
\label{def:convergence1}
    \begin{shaded*}
        A real sequence $(a_n)$ converges to $L\in\R$ if $$\forall\varepsilon>0,\,\exists N\in\N,\, \forall n\ge N,\, |a_n-L|<\varepsilon.$$
    \end{shaded*}
\end{definition}
For notation, if Definition \ref{def:convergence1} holds for a sequence $(a_n)$, we will write $\lim_{n\to\infty}a_n=L$, which we will sometimes shorten to $a_n\to L$ (as $n\to\infty$) or $\lim a_n=L$, if from the context it is clear we are talking about a sequence, and not, say, a limit of a function $f(x)$ (which, as we will see, does not have to be the limit as $x\to\infty$).

\begin{example}
    Use Definition \ref{def:convergence1} to show that the sequence $(a_n)$ defined by $a_n=1/n$ for each $n\in\N$ converges to $0$.
\end{example}
\begin{proof}
    Let $\varepsilon>0$ be arbitary. Choose $N\in\N$ such that $N>1/\varepsilon$ (such an $N$ exists by the Archimedean property of the real numbers from the introductory module).
    Then, given $n\ge N$, we have $$|a_n-0|=1/n\le 1/N<\varepsilon,$$so $a_n\to 0$.\qedhere
\end{proof}

\begin{definition}
    \begin{shaded*}
        A sequence $diverges$ if it does not converge to any $L\in\R$.
    \end{shaded*}
\end{definition}

\begin{definition}
    \begin{shaded*}
        A real sequence $(a_n)$ is bounded if there exists $M>0$ such that for all $n\in \N$, $|a_n|\le M$.
    \end{shaded*}
\end{definition}

\begin{remark}
    Note that a real sequence can diverge but still be bounded (see below).
\end{remark}

\begin{example}
    Use Definition \ref{def:convergence1} to show that the sequence $a_n\coloneqq (-1)^n$ diverges.
\end{example}
\begin{proof}
    Suppose for contradiction that $a_n\to L\in \R$. Let $N\in\N$ be such that for all $n\ge N$, $$\left|(-1)^n-L\right|<1/2.$$
    But also, by the triangle inequality, $$2=\left|(-1)^{2N}-(-1)^{2N+1}\right|\le \left|(-1)^{2N}-L\right|+\left|L-(-1)^{2N+1}\right|<1/2+1/2=1,$$which is a contradiction.
    Therefore, $(a_n)$ diverges.\qedhere
\end{proof}

\subsection{Basic results on convergent sequences}
Unsurprisingly, if a sequence converges to a limit, then it must be unique:
\begin{theorem}
    \begin{shaded*}
        Limits of convergent sequences are unique.
    \end{shaded*}
\end{theorem}
\begin{proof}
    Suppose $a_n\to L$ and $a_n\to M$, and let $\varepsilon>0$ be arbitrary. We wish to conclude that $L=M$. By convergence of $a_n\to L$, pick $N_1\in\N$ such that for all $n\ge N_1$, $|a_n-L|<\varepsilon/2$.
    By convergence of $a_n\to M$ let $N_2\in\N$ be such that for all $n\ge N_2$, $|a_n-M|<\varepsilon/2$. Now, let $n\coloneqq \max\{N_1,N_2\}$. By the triangle inequality, we have
    $$|L-M|=|L-a_n+a_n-M|\le|a_n-L|+|a_n-M|<\frac{\varepsilon}{2}+\frac{\varepsilon}{2}=\varepsilon.$$
    Therefore, since $\varepsilon$ was arbitrary, $L=M$ (to see this, if $L\ne M$ then one can pick $\varepsilon=|L-M|/2$ and get a contradiction).\qedhere
\end{proof}

\begin{theorem}[Order limit theorem]
\label{thm:olt}
    \begin{shaded*}
        Suppose $(a_n)$ and $(b_n)$ are two real sequences such that $a_n\ge b_n$ for all $n\in\N$. Let $a_n\to a$ and $b_n\to b$. Then $a\ge b$.
    \end{shaded*}
\end{theorem}
\begin{proof}
    Suppose for contradiction that $a<b$. Pick $\varepsilon=\frac{b-a}{2}$. Since $a_n\to a$, there exists $N_1\in\N$ such that for all $n\ge N_1$, $|a_n-a|<\varepsilon$, so in particular for $n\ge N_1$, $$a_n<a+\varepsilon=a+\frac{b-a}{2}=\frac{a+b}{2},$$
    and since $b_n\to b$ there exists $N_2\in\N$ such that for all $n\ge N_2$, $|b_n-b|<\varepsilon$, so in particular for $n\ge N_2$, $$b_n>b-\varepsilon=b-\frac{b-a}{2}=\frac{a+b}{2}.$$Therefore, choosing $n\coloneqq\max\{N_1,N_2\}$ it follows that $b_n>a_n$, which is a contradiction.\qedhere
\end{proof}

\begin{corollary}
    \begin{shaded*}
        Suppose $(a_n)$ is a real sequence with $a_n\to a$ such that $B\le a_n\le C$ for all $n\in\N$. Then $B\le a\le C$.
    \end{shaded*}
\end{corollary}
\begin{proof}
    Apply Theorem \ref{thm:olt} using constant sequences $b_n=B$ and $c_n=C$ for all $n\in \N$.\qedhere
\end{proof}
\begin{remark}
    If $B<a_n<C$ for all $n\in\N$ in the above corollary, it is not necessarily true that $B<a<C$, but we do still get the non-strict inequalities $B\le a\le C$. Limits only preserve non-strict inequalities.
\end{remark}

\begin{theorem}
\label{thm:boundedseq}
    \begin{shaded*}
        Convergent sequences are bounded.
    \end{shaded*}
\end{theorem}
\begin{proof}
    Let $(a_n)$ be a convergent sequence, and let $L\in\R$ be such that $\lim a_n = L$. Let $N\in\N$ be such that for all $n\ge N$, $|a_n-L|<1.$ In other words, for all $n\ge N$ we have $$a_n\in (L-1,L+1).$$
    Now, pick $M\coloneqq \max\{|a_1|,|a_2|,\dots,|a_{N-1}|,|L|+1\}$. Then by construction, $|a_n|\le M$ for all $n\in \N$, so $(a_n)$ is bounded.\qedhere
\end{proof}
\begin{remark}
    The idea for the above proof is that we can control $(a_n)$ to be close to $L$ (say, within a distance of $1$ from $L$), past some point $N\in\N$, and before that, there are only finitely many terms in the sequence, so we can just take the maximum of a finite number of terms $|a_1|,\dots,|a_{N-1}|$, and then include $|L|+1$ to complete the bound.
\end{remark}
Similarly to the above observation, changing a finite number of terms in a convergent sequence does not affect its convergence. When dealing with convergence of sequences, we only care about the behaviour of the $tail$ of the sequence.
This motivates the following definition:
\begin{definition}
    \begin{shaded*}
        A sequence $(a_n)$ $eventually$ has a property $P$ if there exists $N\in\N$ such that for all $n\ge N$, $a_n$ has property $P$.
    \end{shaded*}
\end{definition}

\begin{remark}
    Modifying the hypothesis of a theorem which requires a property of a sequence to hold for all $n\in \N$ with $eventually$ does not weaken the theorem; they are equivalent. 
\end{remark}

\begin{example}
    Let the property $P$ be 'constant'. The sequence $a_1=1,a_2=2$ and $a_n=0$ for all $n\ge 3$ is $eventually$ constant.
\end{example}
\begin{exercise*}
    Prove that if an integer sequence $(a_n)$ (that is, $a_n\in\Z$ for all $n\in\N$) converges, then it is eventually constant.
\end{exercise*}



\begin{definition}
\label{def:seqinfinity}
    \begin{shaded*}
        Let $(a_n)$ be a real sequence.
        \begin{itemize}
            \item A sequence $a_n\to\infty$ if for every $R>0$, $a_n$ is eventually greater than $R$.
            \item A sequence $a_n\to-\infty$ if for every $R<0$, $a_n$ is eventually less than $R$.
        \end{itemize}
    \end{shaded*}
\end{definition}

\begin{example}
    \begin{enumerate}
        \item[1.] $a_n=n^2\to\infty$. Indeed, let $R>0$, and take $n>\sqrt{R}$. Then $$a_n=n^2>(\sqrt{R})^2=R$$.
        \item[2.] Define $$a_n=\begin{cases}
            0 & \text{if }n\text{ odd}\\
            n & \text{if }n\text{ even}
        \end{cases}$$
        This is an example of a sequence which does not diverge to $\infty$, but is unbounded (as an exercise, use Definition \ref{def:seqinfinity} to prove this directly).
    \end{enumerate}
\end{example}
\noindent Changing a finite number of terms in a convergent sequence does not affect convergence:

\begin{lemma}
    \begin{shaded*}
        Let $K\in \N$ be arbitrary. Then $(a_n)_{n=1}^{\infty}$ is a convergent sequence with limit $L\in\R$ if and only if $(a_{n+K})_{n=1}^{\infty}$ is a convergent sequence, with limit $L\in\R$.
    \end{shaded*}
\end{lemma}
\begin{proof}
    Suppose $(a_n)$ converges to $L\in\R$. Pick $N\in \N$ such that for all $n\ge N$, $|a_n-L|<\varepsilon$. Since $K\ge 0$, given $n\ge N$, $n+K\ge n \ge N$, so $|a_{n+K}-L|<\varepsilon$.\par
    Conversely, suppose $(a_{n+K})$ converges to $L\in \R$. Let $N\in \N$ be such that for all $n\ge N$, $|a_{n+K}-L|<\varepsilon$. Choose $N'\coloneqq N+K$. Given $n\ge N'$, we have $n\ge N+K$, so for $n\ge N'$, $|a_n-L|<\varepsilon$.\qedhere
\end{proof}




\subsection{Algebraic limit theorem}

There is a useful theorem that allows us to avoid $\varepsilon$ proofs and use convergence results that we have already proven:
\begin{theorem}[Algebraic limit theorem]
\label{thm:alt}
    \begin{shaded*}
        Let $(a_n)$, and $(b_n)$ be real sequences, and suppose $a_n\to a\in\R$, and $b_n\to b\in \R$. Then
        \begin{enumerate}
            \item[1.] $a_n+b_n\to a+b$.
            \item[2.] $a_nb_n\to ab$.
            \item[3.] For all $\lambda\in \R$, $\lambda a_n\to \lambda a$. 
            \item[4.] If $b_n\ne 0$ for all $n\in\N$ and $b\ne 0$, then $1/b_n\to 1/b$. 
        \end{enumerate}
    \end{shaded*}
\end{theorem}
\begin{proof}
    (1) and (3) are good exercises to practice using Definition \ref{def:convergence1}.
    \begin{enumerate}
        \item[(2)] Case 1: without loss of generality, suppose that $a\ne 0$ and $b=0$; let $\varepsilon>0$ be arbitrary. By Theorem \ref{thm:boundedseq}, there exists $M>0$ such that for all $n\in\N$, $|a_n|\le M$.
        By convergence of $b_n\to 0$, let $N\in\N$ be such that for all $n\ge N$, $|b_n|<\varepsilon/M$. Then for $n\ge N$, we have $$|a_nb_n|=|a_n||b_n|\le M|b_n|<M(\varepsilon/M)=\varepsilon,$$
        so $a_nb_n\to 0$.\par
        Case 2: Suppose that $a\ne 0$ and $b\ne 0$; let $\varepsilon>0$ be arbitrary. By Theorem \ref{thm:boundedseq} let $M>0$ be such that for all $n\in\N$, $|a_n|\le M$.
        By convergence of $a_n\to a$, let $N_1\in\N$ be such that for all $n\ge N_1$, $|a_n-a|<\varepsilon/2|b|$.
        Also, by convergence of $b_n\to b$, let $N_2\in\N$ be such that for all $n\ge N_2$, $|b_n-b|<\varepsilon/2M$.
        Choose $N\coloneqq\max\{N_1,N_2\}$. Then for all $n\ge N$, we have
        \begin{align*}
            |a_nb_n-ab|&=|a_nb_n-a_nb+a_nb-ab|\\
            &\le |a_nb_n-a_nb|+|a_nb-ab|\\
            &=|a_n||b_n-b|+|b||a_n-a|\\
            &\le M|b_n-b|+|b||a_n-a|\\
            &< M\left(\frac{\varepsilon}{2M}\right)+|b|\left(\frac{\varepsilon}{2|b|}\right)\\
            &=\varepsilon,
        \end{align*}
        so $a_nb_n\to ab$.
        \item[(4)] Let $\varepsilon>0$ be arbitrary. Let $N_1\in\N$ be such that for all $n\ge N_1$, $|b_n-b|<|b|/2$, so by the triangle inequality we have, for $n\ge N_1$,
        $$|b|=|b-b_n+b_n|\le |b_n-b|+|b_n|<\frac{|b|}{2}+|b_n|\implies |b_n|>\frac{|b|}{2}>0.$$
        Now, let $N_2\in\N$ be such that for all $n\ge N_2$, $|b_n-b|<\frac{\varepsilon|b|^2}{2}$. Choose $N\coloneqq\max\{N_1,N_2\}$. For all $n\ge N$, we have $$\left|\frac{1}{b_n}-\frac{1}{b}\right|=\frac{|b_n-b|}{|b||b_n|}<\frac{|b_n-b|}{|b|^2/2}<\frac{\varepsilon|b|^2/2}{|b|^2/2}=\varepsilon,$$
        so $1/b_n\to 1/b$.\qedhere
    \end{enumerate}
\end{proof}
\begin{remark}
    Using items (2) and (4) above, it follows that if $b\ne 0$ and $b_n\ne 0$ for all $n\in\N$, and $a_n\to a$, $b_n\to b$, then $a_n/b_n\to a/b$.
\end{remark}

\begin{example}
    \begin{enumerate}
        \item[1.] $a_n=\frac{n^2+4n+1}{2n^2+1}\to \frac{1}{2}$. For $n\in\N$ we have $$a_n=\frac{1+4/n+1/n^2}{2+1/n^2}.$$
        Now, using (2) in Theorem \ref{thm:alt}, $1/n^2\to 0$ since we have shown $1/n\to 0$, and using (1), (3), (4) we get $\lim a_n=\frac{1}{2}$.
        \item[2.] Applying (2) inductively, since $1/n\to 0$ we get that for any $m\in\N$, $1/n^m\to 0$.
    \end{enumerate}
\end{example}

\subsection{Monotone convergence theorem}
\begin{definition}
    \begin{shaded*}
        Let $(a_n)$ be a real sequence.
        \begin{enumerate}
            \item[1.] We say that $a_n$ is \textit{increasing} if for all $n\in\N$, $a_{n+1}\ge a_n$. If the inequality is strict for all $n\in\N$, we say that $a_n$ is \textit{strictly increasing}.
            \item[2.] We say that $a_n$ is \textit{decreasing} if for all $n\in\N$, $a_{n+1}\le a_n$. If the inequality is strict for all $n\in\N$, we say that $a_n$ is \textit{strictly decreasing}.
            \item[3.] A sequence $a_n$ is \textit{monotonic} if it is increasing or decreasing. 
        \end{enumerate}
    \end{shaded*}
\end{definition}
\begin{example}
    The sequence $a_n=1/n$ is strictly decreasing and converges to $0$, the sequence $a_n=1-1/n$ is strictly increasing and converges to $1$, the sequence $a_n=n$ is strictly increasing and diverges to $\infty$,
    and the sequence $a_n=(-1)^nn$ is neither strictly increasing nor strictly decreasing, and diverges. Constant sequences are both increasing and decreasing.
\end{example}
\noindent We now state and prove an important theorem about monotonic and bounded real sequences.
\begin{theorem}[Monotone convergence theorem]
\label{thm:mct}
    \begin{shaded*}
        Let $(a_n)$ be a bounded and monotone real sequence. Then $(a_n)$ is convergent.
    \end{shaded*}
\end{theorem}
\begin{proof}
    We will prove the case where $a_n$ is increasing. The case where $a_n$ is decreasing is similar, and is left as an exercise.\par
    Consider the set $S\coloneqq\{a_n:n\in\N\}$. $S$ is nonmempty since $a_1\in S$, and since $a_n$ is bounded, the set $S$ is bounded. By the completeness of $\R$, $s\coloneqq \sup S$ exists. We claim that $a_n\to s$.\par
    Let $\varepsilon>0$. By definition of supremum, there exists $N\in\N$ such that $s-\varepsilon<a_N\le s$. But since $s$ is an upper bound, and $a_n$ is increasing, we therefore have $$s-\varepsilon<a_N\le a_{N+1}\le \cdots\le s<s+\varepsilon,$$
    i.e. for all $n\ge N$, we have $|a_n-s|<\varepsilon$, so $a_n\to s$.\qedhere
\end{proof}

\begin{remark}
    This result is useful because it allows us to deduce convergence of a sequence without knowing what the limit is.
\end{remark}
\begin{example}
    Define a sequence $(a_n)$ by $a_{n+1}=\sqrt{2+a_n}$ for $n\in\N$ with $a_1=\sqrt{2}$. We claim that $(a_n)$ converges.
\end{example}
\begin{proof}
    By induction, $0<a_n<2$ for all $n\in\N$; indeed, if $a_k<2$ for all $k\in\N$ then $a_{k+1}=\sqrt{2+a_k}<\sqrt{2+2}=2$, and the base case $a_1=\sqrt{2}<2$. So $0<a_n<2$ for all $n\in\N$. Therefore, $a_n$ is bounded.
    Also, the sequence is strictly increasing: for $n\in\N$ we have $$a_n=\sqrt{a_na_n}<\sqrt{2a_n}=\sqrt{a_n+a_n}<\sqrt{a_n+2}=a_{n+1},$$
    so by Theorem \ref{thm:mct}, $(a_n)$ is convergent.\par
    In fact, with recurrence relations, we can often find the limit as well, once we know that it converges: this is because $\lim a_{n+1}=\lim a_n=L$ for some $L\in\R$, which must satisfy, by the algebraic limit theorem (Theorem \ref{thm:alt}), $L=\sqrt{2+L}$, i.e. $$L^2-L-2=0\implies (L-2)(L+1)=0,$$
    but since $0<a_n<2$ for all $n\in\N$, $0\le L\le 2$ by Theorem \ref{thm:olt} (order limit theorem), so $L=2$. 
\end{proof}

\subsection{Subsequences}
\begin{definition}
    \begin{shaded*}
        Let $(a_n)$ be a sequence, and let $f:\N\to\N$ be a strictly increasing function. A \textit{subsequence} of $(a_n)$ is a sequence defined by $b_n\coloneqq a_{f(n)}$ for $n\in\N$.
    \end{shaded*}
\end{definition}
\begin{remark}
    In practice, if $\{n_k:k\in\N\}$ is any strictly increasing set of indices such that $n_1<n_2<n_3<\cdots$, the corresponding subsequence of $(a_n)_{n=1}^{\infty}$ is $(a_{n_k})_{k=1}^{\infty}$.
\end{remark}
\begin{example}
    Let $a_n=n$ for $n\in\N$. Then $a_{2n}=2n$ is a subsequence of $a_n$. Let $a_n=1/n$ for $n\in\N$. Then $a_p=1/p$ for $p$ prime is a subsequence of $a_n$. However, $(1/3,1/2,1,1/4,1/5,\dots)$ is not a subsequence of $a_n$ (we changed the ordering of the terms).
\end{example}

\begin{lemma}
\label{lem:subseqlem}
    \begin{shaded*}
        Let $f:\N\to\N$ be a strictly increasing function. Then for all $k\in\N$, $f(k)\ge k$.
    \end{shaded*}
\end{lemma}
\begin{proof}
    The proof is by induction on $k$. For the base case, clearly $f(1)\ge 1$ since the codomain is $\N$, and $1$ is the smallest natural number. Given $m\in\N$, suppose that $f(m)\ge m$.
    Since $f$ is strictly increasing, $f(m+1)>f(m)$ i.e. $f(m+1)\ge f(m)+1$, since $f(m+1)-f(m)\ge 1$ (smallest distance between distinct integers is $1$), so $$f(m+1)\ge f(m)+1\ge m+1,$$by the induction hypothesis, so $f(k)\ge k$ for all $k\in\N$. \qedhere
\end{proof}

\begin{theorem}
\label{thm:subseqlim}
    \begin{shaded*}
        Let $(a_n)$ be a sequence that converges to $L\in\R$. Then any subsequence of $(a_n)$ converges to $L$.
    \end{shaded*}
\end{theorem}
\begin{proof}
    Let $(a_{n_k})$ be an arbitrary subsequence of $(a_n)$, and let $\varepsilon>0$.
    Since $a_n\to L$, pick $N\in\N$ such that for all $n\in\N$, $|a_n-L|<\varepsilon$. Let $k\ge N$. Then by Lemma \ref{lem:subseqlem}, we have $n_k\ge k\ge N$, so $|a_{n_k}-L|<\varepsilon$, so $a_{n_k}\to L$ as $k\to\infty$.\qedhere
\end{proof}
\begin{remark}
    In contrapositive form, Theorem \ref{thm:subseqlim} is very useful for showing that a sequence diverges. You only need to find two subsequences that converge to different limits, to prove that a sequence diverges.
\end{remark}
\begin{example}
    We now get a much easier proof for the divergence of $a_n=(-1)^n$. We have $a_{2n}=1$ for all $n\in\N$ and $a_{2n-1}=-1$ for all $n\in\N$, so by Theorem \ref{thm:subseqlim} (the contrapositive), $a_n$ diverges. 
\end{example}

\begin{theorem}
\label{thm:monotonesubseqthm}
    \begin{shaded*}
        Any sequence has a monotonic subsequence.
    \end{shaded*}
\end{theorem}
\begin{proof}
    Let $(a_n)$ be any sequence. We consider the set of \textit{peak indices} in the sequence, $S\coloneqq\{n\in\N:a_n\ge a_m\, \forall m \ge n\}$. If $k\in S$, we say that $a_k$ is \textit{peak}.\par
    Case 1: $S$ is finite (or empty). Let $m\in\N$ be such that $m>n$ for all $n\in S$ (note that this vacuously includes the case where $S$ is empty). Since $k\notin S$ for all $k\ge m$, in particular for all $k\ge m$, $a_k$ is not \textit{peak}. We therefore construct an increasing subsequence of $(a_n)$ by induction.
    Define $n_1\coloneqq m$. Since $a_{n_1}$ is not \textit{peak}, there exists $n_2>n_1$ such that $a_{n_2}>a_{n_1}$. Given $k\in\N$, suppose that $a_{n_1}<a_{n_2}<\cdots<a_{n_k}$ have been chosen. Since $a_{n_k}$ is not \textit{peak}, there exists $n_{k+1}>n_k$ such that $a_{n_{k+1}}>a_{n_k}$. By construction, this is a (strictly) increasing subsequence of $(a_n)$.\par
    Case 2: $S$ is infinite. By the well-ordering of the natural numbers, we may pick $n_1\coloneqq\min S$. Given $k\in\N$, suppose that $n_1<n_2<\cdots<n_k$ have been chosen such that $n_i\in S$ for all $1\le i\le k$. Since $S$ is infinite, $S\setminus\{n_1,\dots,n_k\}$ is nonempty, so we may pick $n_{k+1}\coloneqq \min (S\setminus\{n_1,\dots,n_k\})$. By definition of the set $S$, for all $k\in\N$, $a_{n_k}\ge a_{n_{k+1}}$, since by construction $a_{n_k}$ is \textit{peak} for all $k\in\N$, so $(a_n)$ is a decreasing subsequence of $(a_n)$.\qedhere
\end{proof}

\begin{theorem}[Bolzano-Weierstraß theorem]
\label{thm:bw}
    \begin{shaded*}
        Any bounded real sequence has a convergent subsequence.
    \end{shaded*}
\end{theorem}
\begin{proof}
    Let $(a_n)$ be a bounded sequence. By Theorem \ref{thm:monotonesubseqthm}, $(a_n)$ has a monotone subsequence $(a_{n_k})$, which is also bounded since $(a_n)$ is bounded, so by the monotone convergence theorem (Theorem \ref{thm:mct}), $(a_{n_k})$ is convergent.\qedhere
\end{proof}
\begin{example}
    The sequence $a_n=\sin (n)$ has a convergent subsequence since $|a_n|\le 1$ for all $n\in\N$. Finding such a subsequence is harder and is beyond the scope of this course.
\end{example}


\subsection{Cauchy sequences}
Convergent sequences, we have seen, get closer and closer to some limit $L$. What can we say about sequences whose terms get arbitrary close together, past some fixed point $N\in\N$? First, we define precisely what this means:
\begin{definition}
    \begin{shaded*}
        A real sequence $(a_n)$ is \textit{Cauchy} if $$\forall\varepsilon>0,\,\exists N\in\N,\, \forall m,n\ge N,\, |a_m-a_n|<\varepsilon.$$
    \end{shaded*}
\end{definition}
A useful fact we can immediately prove (in a similar way to Theorem \ref{thm:boundedseq}) is that Cauchy sequences are bounded.
\begin{lemma}
\label{lem:cauchybounded}
    \begin{shaded*}
        Cauchy sequences are bounded.
    \end{shaded*}
\end{lemma}
\begin{proof}
    Exercise.
\end{proof}

In fact, in $\R$, the concepts of Cauchy sequences and convergent sequences are equivalent in the following sense:
\begin{theorem}[Cauchy criterion for convergence in $\R$]
    \begin{shaded*}
        Let $(a_n)$ be a real sequence. Then $$(a_n)\text{ is convergent}\iff (a_n)\text{ is Cauchy}.$$
    \end{shaded*}
\end{theorem}
\begin{proof}
    Later.
\end{proof}

\begin{remark}
    It is important to note that the above equivalence, in particular Cauchy $\implies$ convergent does not hold for all sets (however, convergence implies Cauchy in any metric space). It relies on the fact that $\R$ is a $complete$ set, that is, it has the least upper bound property.
    For example, $\Q$ is not a complete set (indeed, $\R$ is the completion of $\Q$, as we saw when constructing $\R$ in the introductory module), so the fact that Cauchy sequences are convergent does not hold in $\Q$. For example, $(1+1/n)^n\to e\notin\Q$. We will prove the irrationality of $e$ later, and we may define $e\coloneqq \sup\{(1+1/n)^n:n\in\N\}$; then by the monotonce convergence theorem (Theorem \ref{thm:mct}), $\lim_{n\to\infty}\left(1+\frac{1}{n}\right)^n=e$ (note that there are more practical ways of defining $e$ and the $\exp:\R\to (0,\infty)$ function, as we will see in Chapter 6).
\end{remark}
\begin{remark}
    In addition to the above comment about completeness, we notice that the Bolzano-Weierstraß theorem was required to prove that Cauchy sequences are convergent. In fact, there is an equivalence between the following concepts for complete sets (in particular, $\R$ in our case):
    \begin{itemize}
        \item Least upper bound property: every non-empty subset of $\R$ that is bounded above has a supremum (least upper bound) in $\R$.
        \item Cauchy completeness: Every Cauchy sequence in $\R$ converges to a limit in $\R$.
        \item Bolzano-Weierstraß theorem: Every bounded sequence in $\R$ has a convergent subsequence.
        \item Monotone convergence theorem: Every bounded, monotonic sequence in $\R$ converges.
        \item Nested interval property: Any sequence of nested, closed, and bounded intervals $[a_n,b_n]$ has a nonempty intersection.
    \end{itemize}
\end{remark}
\begin{exercise*}
    Prove that the above concepts are equivalent. 
\end{exercise*}

\subsection{Complex sequences}
As hinted at earlier, sequences need not be real. An interesting and rich area of mathematics is analysis of the complex numbers, $\C$, which we will see more of in second year. 
Here, we build some theory on complex sequences (and later, complex series and complex power series). There will be no rigorous construction of $\C$ (however, it is not too bad now that we have constructed $\R$ because, in fact, we can construct $\C$ from $\R^2$ with some special operations, namely the usual complex addition and complex multiplication). We will assume all the usual properties of the complex numbers, in particular that $\C$ is a field with the usual complex addition and complex multiplication.
\begin{definition}
    \begin{shaded*}
        A complex sequence is a function $f:\N\to\C$.
    \end{shaded*}
\end{definition}
\noindent Again, practically, we write $(z_n)$ or just $z_n$ to denote a complex sequence. For a complex sequence $(z_n)$ it is natural to associate two real sequences $(x_n)$, $(y_n)$ with $(z_n)$ such that for each $n\in \N$, $z_n=x_n+iy_n$, i.e. $x_n=\Re(z_n)$ and $y_n=\Im(z_n)$, where $\Re$ and $\Im$ denote the real and imaginary parts of a complex number, respectively.
It is therefore not too hard to prove facts about complex sequences using our results about real sequences.

\begin{definition}[Convergence of a sequence in $\C$]
    \begin{shaded*}
        A complex sequence $(z_n)$ converges to $z\in \C$ if $$\forall\varepsilon>0,\,\exists N\in\N,\, \forall n\ge N,\, |z_n-z|<\varepsilon.$$
    \end{shaded*}
\end{definition}
\begin{remark}
    Note that the modulus $|\cdot|$ is now the \textit{complex} modulus defined by $|z|\coloneqq\sqrt{x^2+y^2}$ where $z=x+iy$.
\end{remark}

First, we prove a convergence criterion for complex sequences, in relation to its real and imaginary parts:
\begin{theorem}
\label{thm:complexrealconvergence}
    \begin{shaded*}
        Let $(z_n)$ be a complex sequence, and define real sequences $(x_n)$, $(y_n)$ by $x_n\coloneqq \Re(z_n)$ and $y_n\coloneqq\Im(z_n)$ for all $n\in\N$. Let $z\in \C$ and write $z=x+iy$ for some $x,y\in\R$. Then $$z_n\to z\iff x_n\to x\text{ and }y_n\to y.$$
    \end{shaded*}
\end{theorem}
\begin{proof}
    Later.
\end{proof}

\begin{definition}
    \begin{shaded*}
        A complex sequence $(z_n)$ is bounded if there exists $M\in \R_{>0}$ such that $|z_n|\le M$ for all $n\in\N$.
    \end{shaded*}
\end{definition}

\begin{theorem}
    \begin{shaded*}
        Let $(z_n)$ be a complex sequence. If $(z_n)$ converges, then $(z_n)$ is bounded.
    \end{shaded*}
\end{theorem}

\begin{proof}
    Using the parts criterion for convergence in $\C$ (Theorem \ref{thm:complexrealconvergence}), $x_n\coloneqq \Re(z_n)$ and $y_n\coloneqq \Im(z_n)$ are convergent and so by Theorem \ref{thm:boundedseq} $(x_n)$ and $(y_n)$ are bounded. Let $M_1\in \R_{>0}$ be such that $|x_n|\le M_1$ for all $n\in\N$ and let $M_2\in\R_{>0}$ be such that $|y_n|\le M_2$ for all $n\in\N$.
    Choose $M\coloneqq\max\{M_1,M_2\}\in\R_{>0}$; then $|x_n|\le M$ and $|y_n|\le M$ for all $n\in\N$, so for all $n\in\N$ we have $$|z_n|=\sqrt{x_n^2+y_n^2}\le \sqrt{2}M,$$
    so $(z_n)$ is bounded.\qedhere
\end{proof}

\begin{theorem}[Bolzano-Weierstraß theorem in $\C$]
\label{thm:bw-c}
    \begin{shaded*}
        Any bounded complex sequence has a convergent subsequence. 
    \end{shaded*}
\end{theorem}
\begin{proof}
    Let $(z_n)$ be a bounded complex sequence. Put $x_n\coloneqq \Re(z_n)$ and $y_n\coloneqq \Im(z_n)$. Then for all $n\in\N$, $|x_n|=\sqrt{x_n^2}\le\sqrt{x_n^2+y_n^2}=|z_n|$ and $|y_n|\le |z_n|$, so $(x_n)$ and $(y_n)$ are bounded.
    By the Bolzano-Weierstraß theorem in $\R$ (Theorem \ref{thm:bw}), there exists a convergent subsequence $(x_{n_k})$ of $(x_n)$. Consider the subsequence $(y_{n_k})$ of $(y_n)$. By the Bolzano-Weierstraß theorem in $\R$, there exists a convergent subsequence $(y_{n_{k_j}})$ of $(y_{n_k})$.
    Furthermore, by Theorem \ref{thm:subseqlim}, $(x_{n_{k_j}})$ is a convergent subsequence of $(x_n)$, so by the parts criterion for convergence in $\C$, $(z_{n_{k_j}})$ is a convergent subsequence of $(z_n)$.\qedhere
\end{proof}

\begin{definition}
    \begin{shaded*}
        A complex sequence is \textit{Cauchy} if $$\forall\varepsilon>0,\,\exists N\in\N,\,\forall m,n\ge N,\,|z_m-z_n|<\varepsilon.$$
    \end{shaded*}
\end{definition}

\begin{theorem}[Cauchy criterion for convergence in $\C$]
    \begin{shaded*}
        Let $(z_n)$ be a complex sequence and let $x_n\coloneqq \Re(z_n)$, $y_n\coloneqq \Im(z_n)$ for all $n\in\N$. Then the following are equivalent:
        \begin{enumerate}
            \item[1.] $(z_n)$ is Cauchy in $\C$.
            \item[2.] $(x_n)$ and $(y_n)$ are Cauchy in $\R$.
            \item[3.] $(z_n)$ converges in $\C$. 
        \end{enumerate}
    \end{shaded*}
\end{theorem}
\begin{proof}
    Later.
\end{proof}

\section{Series}

\section{Power Series}

\chapter{Limits and Continuity}






\chapter{Differentiable functions}







\chapter{Convexity}







\chapter{The Integral}








\chapter{Taylor's Theorem}









\chapter{Convergence of functions}








\chapter{Series of functions}








\chapter{The Trigonometric functions}


\chapter{Curves in the plane}








\end{document}